\documentclass[../../../include/open-logic-section]{subfiles}

\begin{document}

\olfileid[hi]{sth}{z}{story}
\olsection{कथा का अधिक विस्तृत रूप}
	
\olref[story][approach]{sec} में हमने समुच्चय-निर्माण की प्रक्रिया का
शोएनफ़ील्ड द्वारा दिया वर्णन उद्धृत किया था। अब इस कथा को कुछ अधिक सटीक
बनाने के लिए हम कुछ और सिद्धांत लिखना चाहते हैं। वे ये हैं:
\begin{enumerate}
\item[] \stageshier. प्रत्येक समुच्चय किसी चरण पर बनता है।
\item[] \stagesord. चरण क्रमित हैं: कुछ दूसरे चरणों से \emph{पहले} आते
हैं।\footnote{वास्तव में हम मौन रूप से यह मानेंगे कि चरण \emph{सुव्यवस्थित}
हैं। इसका अर्थ \olref[ordinals][]{chap} में समझाया गया है। यह एक ठोस मान्यता
है। वास्तव में \citet{Scott1974} की अत्यंत चतुर तकनीक से इस मान्यता से
\emph{बचा} जा सकता है और फिर इसे \emph{व्युत्पन्न} किया जा सकता है। (यह यह
भी समझाएगा कि हमें आरंभिक चरण क्यों मानना चाहिए।) यहाँ हम इसका विवेचन नहीं
कर सकते; अधिक के लिए \citet{ButtonLT1} देखें।}
\item[] \stagesacc. किसी भी चरण $S$ और चरण~$S$ से \emph{पहले} बने किन्हीं
भी समुच्चयों के लिए: चरण~$S$ पर ऐसा समुच्चय बनता है जिसके सदस्य ठीक वही
समुच्चय हैं। चरण~$S$ पर और कुछ नहीं बनता।
\end{enumerate}
ये अनौपचारिक सिद्धांत हैं, किंतु इनके उपयोग से हम ज़ेर्मेलो के समुच्चय
सिद्धांत के कई अभिगृहीतों का औचित्य सिद्ध कर सकेंगे।

(यहाँ एक सावधानी आवश्यक है। यद्यपि हम पूरी तरह मानक नामों वाले कुछ पूरी तरह
मानक अभिगृहीत प्रस्तुत करेंगे, अभी प्रस्तुत तिरछे अक्षरों वाले सिद्धांतों के
साहित्य में कोई विशिष्ट नाम नहीं हैं। ये केवल ऐसे उपनाम हैं जिनके उपयोगी
होने की हमें आशा है।)

\end{document}
